\section{The formalized case over \texorpdfstring{$\mathbb F_5$}{F5}}

This section follows the Lean formalization. All polynomials and power series
have coefficients in $\mathbb F_5=\mathbb Z/5\mathbb Z$. The notation is
chosen to agree with the formal definitions.

\subsection{Polynomial orders and column periods}

For $Q\in\mathbb F_5[t]$, define
\[
\polyord(Q)=\polyord_{\mathbb F_5[t]/(Q)}(t),
\]
the multiplicative order of the residue class of $t$. When $Q$ is nonzero and
$Q(0)\ne0$, this is the least positive integer $N$ such that
\begin{equation}\label{eq:f5-order}
Q(t)\mid t^N-1.
\end{equation}
More generally, the formalization proves the divisibility criterion
\begin{equation}\label{eq:f5-order-divisibility}
\polyord(Q)\mid N
\quad\Longleftrightarrow\quad
Q(t)\mid t^N-1.
\end{equation}

Fix
\[
g(t)=\frac{p_1(t)}{p_2(t)},
\qquad
f(t)=\frac{p_3(t)}{p_4(t)},
\]
where $p_2(0)p_4(0)\ne0$. The generating function of column $k$ is
\begin{equation}\label{eq:f5-column}
H_k(t)=\frac{p_1(t)p_3(t)^k}{p_2(t)p_4(t)^k}.
\end{equation}
For an ordinary Riordan array one additionally assumes $g(0)\ne0$,
$f(0)=0$, and $f'(0)\ne0$. The formalized periodicity statements themselves
use only the displayed rational column generating functions and the
nonvanishing denominator constants.
Set
\[
Q_k(t)=p_2(t)p_4(t)^k,
\qquad
\pi_k=\polyord(Q_k).
\]

\begin{prop}\label{prop:f5-eventual}
The coefficient sequence of $H_k$ is eventually $N$-periodic whenever
$\pi_k\mid N$. If
\[
\gcd\bigl(Q_k,p_1p_3^k\bigr)=1,
\]
then every eventual period of $H_k$ is divisible by $\pi_k$. In particular,
$\pi_k$ is its least eventual period.
\end{prop}

The first statement follows by multiplying Equation~\eqref{eq:f5-column} by
$t^N-1$. For the converse, eventual $N$-periodicity makes
$(t^N-1)H_k$ a polynomial. Coprimality then implies
$Q_k\mid t^N-1$, and Equation~\eqref{eq:f5-order-divisibility} gives
$\pi_k\mid N$.

\subsection{The five-adic period tower}

The formalized argument uses
\begin{equation}\label{eq:f5-frobenius}
(t^N-1)^5=t^{5N}-1.
\end{equation}
Since $Q_k\mid Q_{k+1}$, one has $\pi_k\mid\pi_{k+1}$. If $k\geqslant1$,
then
\[
Q_{k+1}=p_2p_4^{k+1}\mid(p_2p_4^k)^2=Q_k^2.
\]
It follows from Equation~\eqref{eq:f5-frobenius} that
$Q_{k+1}\mid t^{5\pi_k}-1$, and hence $\pi_{k+1}\mid5\pi_k$.

\begin{prop}\label{prop:f5-ratio}
Suppose $p_2$ and $p_4$ are nonzero with nonzero constant terms. If
$k\geqslant1$, then
\[
\pi_{k+1}=\pi_k
\qquad\text{or}\qquad
\pi_{k+1}=5\pi_k.
\]
Consequently, there is a nondecreasing function $u\colon\mathbb N\to\mathbb
N$ with $u(1)=0$ such that
\[
\pi_k=\pi_1 5^{u(k)}
\qquad(k\geqslant1).
\]
\end{prop}

The formalization also verifies the factorization formula. Define
\[
\lambda_5(m)=\min\{j\geqslant0:m\leqslant5^j\}
\qquad(m\geqslant1).
\]
If $\varphi$ is irreducible and $\varphi(0)\ne0$, then
\begin{equation}\label{eq:f5-prime-power}
\polyord(\varphi^m)=\polyord(\varphi)5^{\lambda_5(m)}.
\end{equation}
Suppose a common finite family of pairwise nonassociated irreducibles gives
\[
p_2=\prod_{i\in I}\varphi_i^{a_i},
\qquad
p_4=\prod_{i\in I}\varphi_i^{b_i},
\]
with $\varphi_i(0)\ne0$ and $a_i+kb_i\geqslant1$. Then
\begin{align}
\pi_k
&=\mathop{\operatorname{lcm}}_{i\in I}
  \left(\polyord(\varphi_i)5^{\lambda_5(a_i+kb_i)}\right)\label{eq:f5-lcm}\\
&=E5^{s_k},\label{eq:f5-explicit}
\end{align}
where
\[
E=\mathop{\operatorname{lcm}}_{i\in I}\polyord(\varphi_i),
\qquad
s_k=\max_{i\in I}\lambda_5(a_i+kb_i).
\]
The second equality uses the fact that each $\polyord(\varphi_i)$ is relatively
prime to five.

\subsection{Circulant blocks}

Let a power series $H(t)=\sum_{n\geqslant0}c_nt^n$ be eventually
$N$-periodic. For a starting index $M$ after the periodic tail has begun,
define
\[
B_{N,M}(H)=\sum_{r=0}^{N-1}c_{M+r}t^r
\quad\text{in}\quad
\mathcal C_N=\mathbb F_5[t]/(t^N-1).
\]
The formalization proves that taking a sufficiently late block intertwines
multiplication by a polynomial with multiplication by its class in
$\mathcal C_N$. It also proves that moving the starting index corresponds to
cyclic shift.

The column generating functions satisfy
\begin{equation}\label{eq:f5-neighbors}
p_4(t)H_{k+1}(t)=p_3(t)H_k(t).
\end{equation}

\begin{prop}\label{prop:f5-block}
Let $N=\pi_{k+1}$. There is an index $M_0$ such that, for every
$M\geqslant M_0$, the length-$N$ blocks read from columns $k$ and $k+1$
satisfy
\[
p_4(t)B_{N,M}(H_{k+1})=p_3(t)B_{N,M}(H_k)
\quad\text{in }\mathcal C_N.
\]
Moreover, $B_{N,M}(H_k)$ is the length-$\pi_k$ block repeated
$N/\pi_k$ times.
\end{prop}

If $p_4$ is nonconstant, its class in $\mathcal C_N$ is not a unit because
$p_4\mid Q_{k+1}\mid t^N-1$. Thus the displayed relation cannot generally be
solved by inverting $p_4$.

\subsection{Pascal's array modulo five}

For
\[
g(t)=\frac{1}{1-t},
\qquad
f(t)=\frac{t}{1-t},
\]
one has
\[
H_k(t)=\frac{t^k}{(1-t)^{k+1}},
\qquad
\pi_k=5^{\lambda_5(k+1)}.
\]
The first periods are
\[
1,5,5,5,5,25,\ldots,25,125,\ldots.
\]
For $N=\pi_{k+1}$, sufficiently late neighboring blocks satisfy
\[
(1-t)B_{k+1}=tB_k
\quad\text{in }\mathbb F_5[t]/(t^N-1).
\]
This is a cyclic discrete-antidifference equation. The formalization also
checks directly that the class of $1-t$ is not a unit in the corresponding
circulant ring.
